% !PRINTED 142
\weberpage{142}
%<<CONT
the two invariants $p, p$ or with one invariant $p^2$), and,
if $p$ is different from $q$, one group (with the two invariants
$p, q$).

There thus remains for us the case
\[
p \equiv 1\ (\mathrm{mod}\ q)
\]
still to investigate, in which $\nu$ can have any value satisfying the condition
\[
\nu \equiv 0\ \left(\mathrm{mod}\ \frac{p-1}{q}\right)
\]
of which there are $q$ different ones with respect to the modulus
$p - 1$, namely:
\weq{7}{\nu = 0,\quad \frac{p-1}{q},\quad 2\,\frac{p-1}{q},\quad \dots\ \frac{(q-1)(p-1)}{q}.}
The case $\nu = 0$ leads here too to an \gesperrt{Abelian} group.
But the other values of $\nu$ are also admissible, and they lead
to just as many non-commutative groups. These non-commutative
groups are isomorphic to one another, and are reduced
to one another when one replaces $b$ by a $b^\beta$,
as the formula (5) shows.

We may assume $b^q = 1$, for $b^q$ is certainly contained in $Q$,
and $b^{pq} = 1$. If, then, $b^q$ is not $= 1$, we replace
$b$ by $b^p$.

One obtains one of the non-commutative groups by letting,
in
\weq{8}{\Theta = a^\alpha b^\beta}
$\alpha$ run through the sequence of numbers $0, 1, \dots p - 1$ and $\beta$ through the numbers
$0, 1, \dots q - 1$. The composition of two
elements of this group results, by (5), from
\weq{9}{b^\beta a^\alpha = a^{\alpha g^{-\nu\beta}} b^\beta, \quad a^p = 1,\ b^q = 1,}
by which every compositum of elements $\Theta$ can be reduced to the
form $\Theta$.

To show that the elements (8) really form a group under
the rules of composition (9), one has to prove the properties
of the group, namely that from $\Theta\Theta' = \Theta\Theta''$ and from $\Theta'\Theta = \Theta''\Theta$
it follows that $\Theta' = \Theta''$, and further the associative law
\[
\Theta(\Theta'\Theta'') = (\Theta\Theta')\Theta''
\]
Both, however, follow very easily from the composition
that follows from (9):
\weq{10}{a^\alpha b^\beta a^{\alpha'} b^{\beta'} = a^{\alpha+\alpha' g^{-\nu\beta}} b^{\beta+\beta'}.}
